\documentclass[10pt,a4paper]{article}
\usepackage{amsmath,amssymb,amsthm,graphicx}
\usepackage{amsthm}
\usepackage{amsfonts, amsmath, amssymb, amsthm}
\usepackage{hyperref}
\usepackage{cleveref}
\def\TT{\mathbb{T}}
\def\PP{\mathbb{P}}
\def\DD{\mathbb{D}}
\def\QQ{\mathbb{Q}}
\def\WW{\mathbb{W}}
\def\MM{\mathbb{M}}
\def\HH{\mathbb{H}}

\newcommand{\Psf}{\mathsf{P}}
\newcommand{\Qsf}{\mathsf{Q}}
\newcommand{\Rsf}{\mathsf{R}}


\def\F{\mathcal{F}}
\def\U{\mathcal{U}}
\def\C{\mathcal{C}}
\def\I{\mathcal{I}}
\def\D{\mathcal{D}}
\def\L{\mathcal{L}}
\def\V{\mathcal{V}}
\def\A{\mathcal{A}}
\def\B{\mathcal{B}}
\def\P{\mathcal{P}}
\def\M{\mathcal{M}}
\def\Nc{\mathcal{N}}
\def\Q{\mathcal{Q}}

% Baire space, Cantor space, binary strings
\newcommand{\cs}{2^{\NN}}
\newcommand{\bs}{\NN^{\NN}}
\newcommand{\bstr}{2^{<\NN}}

\newcommand{\uh}{\upharpoonright}

% Turing machines
\newcommand{\converge}{\!\!\downarrow}
\newcommand{\diverge}{\!\!\uparrow}

% Forcing question
\newcommand{\qvdash}{\operatorname{{?}{\vdash}}}
\newcommand{\nqvdash}{\operatorname{{?}{\nvdash}}}

\newcommand{\ISig}{\mathsf{I}\Sigma^0}
\newcommand{\Robinson}{\mathsf{Q}}
\newcommand{\IPi}{\mathsf{I}\Pi^0}
\newcommand{\IDelta}{\mathsf{I}\Delta^0}
\newcommand{\Iopen}{\mathsf{I_{open}}}
\newcommand{\BSig}{\mathsf{B}\Sigma^0}
\newcommand{\BPi}{\mathsf{B}\Pi^0}
\newcommand{\BDelta}{\mathsf{B}\Delta^0}
\newcommand{\BCSig}{\mathsf{BC}\Sigma^0}
\newcommand{\BCPi}{\mathsf{BC}\Pi^0}
\newcommand{\BCDelta}{\mathsf{BC}\Delta^0}
\newcommand{\LSig}{\mathsf{L}\Sigma^0}
\newcommand{\LPi}{\mathsf{L}\Pi^0}
\newcommand{\CSig}{\mathsf{C}\Sigma}

\newcommand{\RCA}[0]{\mathsf{RCA}}
\newcommand{\WKL}[0]{\mathsf{WKL}}
\newcommand{\WWKL}[0]{\mathsf{WWKL}}
\newcommand{\ACA}[0]{\mathsf{ACA}}
\newcommand{\ATR}[0]{\mathsf{ATR}}
\newcommand{\PIOOCA}[0]{\Pi^1_1\mathsf{-CA}}
\newcommand{\CAC}[0]{\mathsf{CAC}}
\newcommand{\ADS}[0]{\mathsf{ADS}}
\newcommand{\RRT}[0]{\mathsf{RRT}}
\newcommand{\FS}[0]{\mathsf{FS}}
\newcommand{\PFS}[0]{\mathsf{PFS}}
\newcommand{\TS}[0]{\mathsf{TS}}
\newcommand{\GRRT}[0]{\mathsf{GRRT}}
\newcommand{\LRRT}[0]{\mathsf{LRRT}}
\newcommand{\EM}[0]{\mathsf{EM}}
\newcommand{\SEM}[0]{\mathsf{SEM}}
\newcommand{\HalfEM}[0]{\mathsf{Half}\mbox{-}\mathsf{EM}}
\newcommand{\RT}[0]{\mathsf{RT}}
\newcommand{\SRT}[0]{\mathsf{SRT}}
%\newcommand{\RRT}{\mathsf{RRT}}

\newcommand{\NN}[0]{\mathbb{N}}
\newcommand{\nat}[0]{\mathbb{N}}

\newcommand{\card}{\operatorname{card}}
\newcommand{\finsub}{\subseteq_{\mathtt{fin}}}

% reductions

\newcommand{\W}{\mathrm{W}}
\newcommand{\sW}{\mathrm{sW}}

\newcommand{\lex}{\mathrm{lex}}

% Barriers
\newcommand{\base}{\operatorname{base}}
\newcommand{\range}{\operatorname{ran}}

%typographically correct double quotes
\def\qt#1{``#1''}%
\newtheorem*{statement}{Statement}
\newtheorem{theorem}{Theorem}
\numberwithin{theorem}{section}
\newtheorem{maintheorem}[theorem]{Main Theorem}
\newtheorem{lemma}[theorem]{Lemma}
\newtheorem{question}[theorem]{Question}
\newtheorem{proposition}[theorem]{Proposition}
%\newtheorem{statement}[theorem]{Statement}
\newtheorem{remark}[theorem]{Remark}
\newtheorem{definition}[theorem]{Definition}
\newtheorem{corollary}[theorem]{Corollary}
\newtheorem{example}[theorem]{Example}

\newcommand{\R}{\mathcal{R}}
\renewcommand{\S}{\mathcal{S}}
\newcommand{\T}{\mathcal{T}}
\title{Strong cone avoidance for computably bounded progressive colorings of variable-sized finite sets}
\author{Lorenzo Carlucci, Oriola Gjetaj, Quentin Le Houérou, Ludovic Levy Patey}
\date{arxiv:2412.11598v3; CC BY 4.0}
\begin{document}
\maketitle
\noindent\textit{Editorial note.} The coloring f need not be Z-computable. The bounding function h is Z-computable; this distinction is essential to strong cone avoidance.\par\medskip
\noindent\textit{Editorial note.} All new forcing conditions, stabilization/extension lemmas, relation coding, specification and diagonalization arguments are included in the solution. No forcing lemma is supplied as a free prerequisite.\par\medskip
\noindent\textit{Editorial note.} The target is the progressive-coloring theorem itself. Its rainbow corollary and finite-dimensional special cases are not counted separately.\par\medskip
\noindent\textit{Editorial scope.} One original result. Full source-local mathematical construction included; named classical inputs are explicit.\par\medskip
\noindent\textit{Allowed standard prerequisites (editorial).}\begin{itemize}
\item Finite-dimensional strong cone avoidance. Wang strong cone avoidance for fixed-arity free sets; Cholak-Patey strong cone avoidance for the achromatic Ramsey theorem at the Catalan bound C\_n. Both cited classical statements are supplied verbatim.
\item Effective compactness and countable forcing. König compactness for finitely branching trees, Sigma01(X) enumeration, and constructing a descending chain meeting countably many dense requirements.
\end{itemize}\medskip
There exist four other subsystems which, together with $\RCA_0$, calibrate the strength of most theorems. These subsystems are known as the \qt{Big Five} (see Mont\'alban~\cite{montalban_open_2011}). Among these, $\ACA_0$ (Arithmetic Comprehension Axiom) extends the axioms of $\RCA_0$ with the comprehension scheme for all arithmetic formulas. Based on the correspondence between computability and definability, it is equivalent to stating the existence of the Turing jump of any set ($\forall X \exists Y (Y = X')$).

We shall also consider two lesser-known stronger variants of $\ACA_0$, namely, $\ACA_0'$ and $\ACA_0^+$. The system $\ACA_0'$ extends $\RCA_0$ with the axiom stating closure under all finite jumps; the system $\ACA_0^+$ extends $\RCA_0$ with the axiom stating the existence of the $\omega$-jump of any set. The $\omega$-jump of a set~$X$ is the set $X^{(\omega)} = \bigoplus_n X^{(n)}$. The system $\ACA_0^+$ is famous for being the best known upper bound to Hindman's theorem.%\lorenzo{I added mention of ACA0' since we mention it at least once.}

Let us first fix some notation and recall the definition of these principles for colorings of finite subsets of fixed size. We denote by $\NN$ the set of natural numbers and by $\NN^+$ the set of positive natural numbers. For $X\subseteq \NN$ and $n \in \NN$
we denote by $[X]^{n}$ the set of all $n$-subsets of $X$. We denote by $[X]^{<\omega}$ the set of all finite subsets of $X$. We always assume that sets are presented in increasing order and make no distinction between a set and the sequence of its elements in increasing order. We write $x \cdot y$ for the concatenation of two sequences $x$ and $y$. We identify a natural number with the set of its predecessors, so that, if $k\geq 1$ we can write $f: X \to k$ to declare a function from $X$ to $\{0, 1, \dots, k-1\}$. A function of this type is often called a coloring of $X$ in $k$ colors.

Given a set~$X$ and $n \in \NN$, we write $[X]^{\leq n}$ for the set of all subsets $s \subseteq X$ such that $|s| \leq n$, and $[X]^{\leq !\omega}$ for the set of all subsets $s \subseteq X$ such that $|s| \leq 1+\min s$. By convention, for~$n = 0$, $[X]^{\leq n}$ is the singleton $\{\emptyset\}$. Given a function $h : \NN \to \NN$, we write $[X]^{\leq h(\cdot)}$ for the set of all finite $s \subseteq X$ such that $|s| \leq h(\min s)$. %\lorenzo{why don't we just write $[X]^{\leq n}$ for $[X]^{!\leq n}$? Same for $\leq h$}\ludovic{This is legacy from when I was using $\alpha \leq \omega$. we can remove the exclamation point. I will do it if you want.}
%\lorenzo{I find notation $[X]^{\leq h(\cdot)}$ a bit confusing. Is $[X]^{\leq h(\cdot)}$ viable for you?}\ludovic{Yes, I will change it}

\begin{definition}
A set $B \subseteq [\NN]^{<\omega}$ is \emph{$\omega$-bounded} if and only if $B \subseteq [\NN]^{\leq h(\cdot)}$ for some function~$h : \NN \to \NN$. It is \emph{computably $\omega$-bounded} if furthermore $h$ is computable.
\end{definition}

\begin{definition}
A problem~$\Psf$ admits \emph{strong cone avoidance} if for every set~$Z$, every non-$Z$-computable set~$C$ and every $\Psf$-instance~$X$, there exists a $\Psf$-solution~$Y$ to~$X$ such that $C \not \leq_T Y \oplus Z$.
\end{definition}

Note that in the previous definition, no computability constraint is given on the $\Psf$-instance~$X$. Thus, strong cone avoidance reflects the combinatorial weakness of~$\Psf$, in the sense that no matter how complex the instance is, it cannot code in its solutions an infinite binary sequence. We shall use the following two theorems:

\begin{theorem}[Wang~\cite{wang2014some}]
For every~$n \in \NN$, $\FS^n$, $\TS^n$ and $\RRT^n$ admit strong cone avoidance.
\end{theorem}

\begin{definition}[Achromatic Ramsey Theorem]
Let $n, k, \ell \in \NN^+$. For every coloring $f : [\NN]^n \to k$, there exists an infinite $H \subseteq \NN$
such that $|f([H]^n)| \leq \ell$. We denote this statement $\RT^n_{k, \ell}$. We write $\RT^n_{<\infty, \ell}$ for $\forall k \RT^n_{k, \ell}$.
\end{definition}

The following sequence of numbers, known as \emph{Catalan numbers}, is omnipresent in Combinatorics. It is inductively defined as follows:
$$
C_0 = 1 \hspace{1cm} C_{n+1} = \sum_{i = 0}^n C_i C_{n-i}
$$
This sequence starts with $1, 1, 2, 5, 14, 42, \dots$ (see sequence A000108 in the OEIS). The number $C_n$ admits many characterizations, such as the number of ways of associating $n$ applications of a binary operator. In computability theory, the $n$th Catalan number $C_n$ surprisingly arose as the exact threshold $\ell$ at which $\RT^n_{<\infty, \ell}$ admits strong cone avoidance.

\begin{theorem}[Cholak and Patey~\cite{cholak2020thin}]
For every~$n \in \NN^+$, $\RT^n_{<\infty, C_n}$ admits strong cone avoidance.
\end{theorem}

\begin{definition}
Let $B \subseteq [\NN]^{<\omega}$ be a set and $f : B \to [\NN]^{<\omega}$ be a coloring.
A set~$H \subseteq \base(B)$ is \emph{$f$-free} if for every~$s \in B \cap [H]^{<\omega}$, $f(s) \cap H \subseteq s$.
\end{definition}

For example, given a coloring $f : [\NN]^n \to \NN$, one can let $B = [\NN]^n$ and $g : [\NN]^n \to [\NN]^{<\omega}$ be defined by $g(s) = \{f(s)\}$. Then a set is $f$-free if and only if it is $g$-free.
Of course, even with $B = [\NN]^n$, infinite free sets do not necessarily exist for arbitrary colorings. We need to impose some constraints on the size of the sets in the image of~$f$.

\begin{definition}
Fix $B \subseteq [\NN]^{<\omega}$.
A coloring~$f : B \to [\NN]^{<\omega}$ is \emph{$b$-constrained} for a bounding function $b : \NN \to \NN$ if for every~$s \in B$, $|f(s)| \leq b(\min s)$. If $b$ is the constant function $x \mapsto k$, then we say that $f$ is \emph{$k$-constrained}.
\end{definition}



\begin{definition}[$k$-Constrained Free Set Theorem]
Fix $B \subseteq [\NN]^{<\omega}$ and $k \in \NN$.
$\FS^B_k$ is the statement \qt{For every coloring  $f : B \to [\NN]^{\leq k}$, there is an infinite $f$-free set}.
\end{definition}

Given $n \in \NN$ and a function $h : \NN \to \NN$, we write $\FS^{\leq n}_k$, $\FS^{\leq !\omega}_k$ and $\FS^{\leq h(\cdot)}_k$ for $\FS^B_k$ when $B$ is $[\NN]^{\leq n}$, $[\NN]^{\leq !\omega}$ and $[\NN]^{\leq h(\cdot)}$, respectively.

\begin{definition}
Fix $B \subseteq [\NN]^{<\omega}$.
A function $f : B \to [\NN]^{<\omega}$ is \emph{progressive} if for every~$s \in B$, either $f(s) = \emptyset$, or $\min f(s)\geq \min s$.
\end{definition}

An easy combinatorial argument shows that the following statement is classically true.


\begin{definition}[Progressive Free Set Theorem]
Fix a set $B \subseteq [\NN]^{<\omega}$ and a bounding function $b : \NN \to \NN$.
$\PFS^B_b$ is the statement \qt{For every $b$-constrained progressive coloring $f : B \to [\NN]^{<\omega}$, there is an infinite $f$-free set}.
\end{definition}

\begin{theorem}\label[theorem]{thm:pfs-omega-strong-cone-avoidance}
Fix a set~$Z$, a non-$Z$-computable set~$D$ and a $Z$-computable function $h : \NN \to \NN$.
For every $h$-constrained progressive coloring $f : [\NN]^{\leq h(\cdot)} \to [\NN]^{<\omega}$, there exists an infinite $f$-free set $G \subseteq \NN$ such that $D \not \leq_T G \oplus Z$.
\end{theorem}

\begin{proposition}\label[proposition]{prop:fs-n-k-strong-cone-avoidance}
For every $k, n \in \NN$, $\FS^{\leq n}_k$ admits strong cone avoidance.
\end{proposition}
\begin{proof}
Fix some set~$Z$, some non-$Z$-computable set~$C$, and some coloring $f : [\NN]^{\leq n} \to [\NN]^{\leq k}$.
For every~$m \leq n$ and $j < k$, let $f_{m, j} : [\NN]^m \to \NN$ be the coloring defined for every~$s \in [\NN]^m$ by letting $f_{m, j}(s)$ be the $j$th element of~$f(s)$, if it exists, and $f_{m, j}(s) = 0$ otherwise. By finitely many successive applications of strong cone avoidance of $\FS^m$ for $m \leq n$ (see Wang~\cite{wang2014some}), there is an infinite set~$H \subseteq \NN$ which is simultaneously $f_{m,j}$-free for every $m \leq n$ and $j < k$, and such that $C \not \leq_T H \oplus Z$.

We claim that $H$ is $f$-free. Suppose by way of contradiction that there is some~$s \in [H]^{\leq n}$ and some~$c \in (f(s) \cap H) \setminus s$. %Note that $s$ is necessarily non-empty. 
Let $j$ be such that $c$ is the $j$th element of~$f(s)$, and let $m = |s|$. Then $f_{m,j}(s) = c$, contradicting $f_{m,j}$-freeness of~$H$.
\end{proof}

\begin{proof}
For simplicity, we prove the theorem in a non-relativized form. Relativization is straightforward.
Fix a non-computable set~$D$. 
Let $\P$ be the collection of all progressive colorings of type $[\NN]^{\leq h(\cdot)} \to [\NN]^{<\omega}$. For a function $b:\NN\to\NN$ we write $\P_b$ for the class of all $b$-constrained colorings in~$\P$.

Given two colorings $f, g \in \P$, we write $g \leq f$ if for every~$s \in [\NN]^{\leq h(\cdot)}$, $g(s) \supseteq f(s)$.
Note that if $H$ is $g$-free and $g \leq f$, then $H$ is $f$-free.
Given two colorings $f, g \in \P$, let $f \cup g$ be the coloring defined by $(f \cup g)(s) = f(s) \cup g(s)$. The coloring $f \cup g$ is the greatest lower bound of~$f$ and $g$ with respect to $\leq$.


Consider the following notion of forcing:
%and an arbitrary coloring $f \in \P_b$.
%We shall construct an infinite $f$-free set~$G \subseteq \omega$ such that $C \not \leq_T G$ using the following notion of forcing:

\begin{definition}
A \emph{condition} is a triple $(f, \sigma, X)$ such that
$f \in \P$ is $b_f$-constrained, for some computable function $b_f : \NN \to \NN$,
$\sigma \in [\NN]^{<\omega}$, $X \subseteq \NN$ is an infinite set such that $\max \sigma < \min X$, and
\begin{enumerate}
    \item[(a)] for every~$s \in [\sigma \cup X]^{\leq h(\cdot)}$ with $\min s \in \sigma$, $f(s) \cap X \subseteq s$.
    \item[(b)] for every~$s \in [\sigma \cup X]^{\leq h(\cdot)}$, $f(s) \cap \sigma \subseteq s$.
    \item[(c)] $D \not \leq_T X$.
\end{enumerate}
A condition $d = (g, \tau, Y)$ \emph{extends} $c = (f, \sigma, X)$ (written $d \leq c$) if
$g \leq f$, $\tau \succeq \sigma$, $Y \subseteq X$ and $\tau \setminus \sigma \subseteq X$.
\end{definition}

The following lemma states that property (a) can be obtained \qt{for free}, that is, by restricting the reservoir, and therefore does not impose any constraint on the stem.
In what follows, fix a coloring $f \in \P$ which is $b_f$-constrained, for some computable function $b_f : \NN \to \NN$.

\begin{lemma}\label[lemma]{lem:fs-preserving-a}
For every $\sigma \in [\NN]^{<\omega}$ and every infinite set~$X$ such that $\max \sigma < \min X$ and $D \not \leq_T X$,
there is an infinite set $Y \subseteq X$ such that $D \not \leq_T Y$ and 
for every~$s \in [\sigma \cup Y]^{\leq h(\cdot)}$ with $\min s \in \sigma$, $f(s) \cap Y \subseteq s$.
\end{lemma}
\begin{proof}
For every~$t \in [\sigma]^{\leq h(\cdot)}$ with $t \neq \emptyset$, let $f_t : [X]^{\leq h(\min t) - |t|} \to [\NN]^{\leq b_f(\min t)}$ be defined for every~$u$ by $f_t(u) = f(t \cdot u)$.
By finitely many successive applications of \Cref{prop:fs-n-k-strong-cone-avoidance},
there exists an infinite subset~$Y \subseteq X$ such that $D \not \leq_T Y$
and such that $Y$ is simultaneously $f_t$-free for every~$t \in [\sigma]^{\leq h(\cdot)}$ with $t \neq \emptyset$.

We claim that $Y$ is our desired set. Indeed, for every~$s \in [\sigma \cup Y]^{\leq h(\cdot)}$ with $\min s \in \sigma$, letting $t = s \cap \sigma$ and $u = s \cap Y$, we have $f(s) = f_t(u)$, so by $f_t$-freeness of~$Y$, $f(s) = f_t(u) \cap Y \subseteq u \subseteq s$.
\end{proof}

In what follows, recall that $C_n$ stands for the $n$th Catalan number.

%\lorenzo{there was a spare $\mu$ in the below definition probably from some previous version, which I changed to $t$}
\begin{definition}
A set $X$ \emph{stabilizes} $\sigma$ if for every $t \in [\sigma]^{\leq h(\cdot)}$ with $t \neq \emptyset$ and every~$n \leq h(\min t) - |t|$,
there is a set $I_{t,n} \subseteq [\sigma]^{\leq b_f(\min t) \times C_n}$ such that for every~$u \in [X]^n$,
$f(t, u) \cap \sigma \subseteq I_{t, n}$.
\end{definition}

\begin{lemma}\label[lemma]{lem:fs-stabilizing}
For every $\sigma \in [\NN]^{<\omega}$ and every infinite set~$X$ such that $\max \sigma < \min X$ and $D \not \leq_T X$,
there is an infinite set $Y \subseteq X$ stabilizing $\sigma$ such that $D \not \leq_T Y$.
\end{lemma}
\begin{proof}
For every~$t \in [\sigma]^{\leq h(\cdot)}$ with $t \neq \emptyset$ and every~$n \leq h(\min t) - |t|$, let $g_{t, n} : [X]^n \to [\sigma]^{\leq b_f(\min t)}$ be defined by $g_{t, n}(u) = f(t \cdot u) \cap \sigma$. By finitely many successive applications of strong cone avoidance of $\RT^n_{<\infty, C_n}$ (see Cholak and Patey~\cite{cholak2020thin}), there is an infinite subset $Y \subseteq X$ such that $D \not \leq_T Y$ and for every~$t \in [\sigma]^{\leq h(\cdot)}$ with $t \neq \emptyset$ and every~$n \leq h(\min t) - |t|$, $|g_{t, n}([Y]^n)| \leq C_n$. Note that here, a color is an element of $[\sigma]^{\leq b_f(\min t)}$ instead of a natural number. Let $I_{t, n} = \bigcup g_{t, n}([Y]^n)$. Then $|I_{t, n}| \leq b_f(\min t) \times C_n$. By definition, for every~$u \in [Y]^n$, $f(t \cdot u) \cap \sigma = g_{t, n}(u) \subseteq I_{t, n}$.
\end{proof}

Given a computable function $b : \NN \to \NN$, let $b ^{+} : \NN \to \NN$ be the computable function defined by  $b^{+}(m) = \sum_{n \leq h(m)} b(m) \times C_n$.

\begin{definition}
Let $X$ be a reservoir stabilizing $[0, k]$. The \emph{limit coloring} is the function $g_{k,X} : [k]^{\leq h(\cdot)} \to [k]^{<\omega}$ defined by $g_{k,X}(\emptyset) = \emptyset$ and $g_{k,X}(t) = \bigcup_{n \leq h(t) - |t|} I_{t, n}$ otherwise.
\end{definition}

The limit function $g_{k,X}$ is $b_f^{+}$-constrained. Note that if $\rho \subseteq [0, k]$ is $g_{k,X}$-free, then it is $f$-free. Indeed, for every~$t \in [0, k]^{\leq h(\cdot)}$, $g_{k,X}(t) = I_{t, 0}$ and by definition of stability for~$n = 0$, $f(t) \cap [0, k] \subseteq I_{t, 0}$.

The following lemma is the core combinatorial lemma which specifies the conditions under which a block of elements~$\rho$ can be added to the stem while preserving the property~(b).

\begin{lemma}\label[lemma]{lem:fs-preserving-b}
Let $(f, \sigma, X)$ be a condition
and $Y \subseteq X$ be an infinite set stabilizing $[0, k]$ for some~$k \in \NN$ and let $g_{k,Y} : [k]^{\leq h(\cdot)} \to [k]^{<\omega}$ be the limit coloring.
Let $\rho \subseteq X \uh_k$ be a finite $g_{k,Y}$-free set. Then $(f, \sigma \cup \rho, Y)$
satisfies property~(b).
\end{lemma}
\begin{proof}
Fix some~$s \in [\sigma \cup \rho \cup Y]^{\leq h(\cdot)}$. We have multiple cases.
\begin{itemize}
    \item Case 1: $s \cap \sigma \neq \emptyset$. Then by properties (a) and (b) of $(f, \sigma, X)$, $f(s) \cap (\sigma \cup X) \subseteq s$. Since $(\sigma \cup \rho) \subseteq (\sigma \cup X)$, then $f(s) \cap (\sigma \cup \rho) \subseteq s$.
    \item Case 2: $s \cap \sigma = \emptyset$ but $s \cap \rho \neq \emptyset$. By property (b) of $(f, \sigma, X)$, $f(s) \cap \sigma \subseteq s$. Let $t = s \cap \rho$ and $u = s \cap Y$. Since $\rho$ is $g_{k,Y}$-free, $g_{k,Y}(t) \cap \rho \subseteq t$. By definition of $g_{k,Y}(t)$, $I_{t, |u|} \subseteq g_{k,Y}(t)$, so $I_{t, |u|} \cap \rho \subseteq t$. In particular, $f(s) = f(t, u) \in I_{t, |u|}$, so $f(s) \cap \rho \subseteq t \subseteq s$. Thus, $f(s) \cap (\sigma \cup \rho) \subseteq s$.
    \item Case 3: $s \cap (\sigma \cup \rho) = \emptyset$. Then $s \subseteq Y$. Since $\min Y > \max (\sigma \cup \rho)$, then by progressiveness of~$f$, $f(s) \cap (\sigma \cup \rho) = \emptyset$.
\end{itemize}
\end{proof}

One can combine \Cref{lem:fs-preserving-a,lem:fs-stabilizing,lem:fs-preserving-b} to obtain an extensibility lemma, saying that every sufficiently generic filter induces an infinite set. 

\begin{lemma}\label[lemma]{lem:fs-extendibility}
Let $(f, \sigma, X)$ be a condition. There is an extension $(f, \tau, Y) \leq (f, \sigma, X)$ such that $|\tau| > |\sigma|$.
\end{lemma}
\begin{proof}
Let $x = \min X$. By \Cref{lem:fs-stabilizing}, there is an infinite subset $Y_0 \subseteq X$ stabilizing $[0, x]$ such that $D \not \leq_T Y_0$. Let $g$ be the limit function. By \Cref{lem:fs-preserving-a}, there is an infinite subset $Y \subseteq Y_0$ such that $(f, \sigma \cup \{x\}, Y)$ satisfies property (a).
By \Cref{lem:fs-preserving-b}, $\{x\}$ being vacuously $g$-free, $(f, \sigma \cup \{x\}, Y)$ satisfies property (b). Thus, $(f, \sigma \cup \{x\}, Y)$ is a valid extension.
\end{proof}

Given a computable function $b : \NN \to \NN$, the space $\P_b$ is not compact.
However, $\P_b$ is in one-to-one correspondence with the effectively compact space $\R_b$ of all relations $R \subseteq [\NN]^{\leq h(\cdot)} \times \NN$ such that for every~$s \in [\NN]^{\leq h(\cdot)}$, $|\{ y : (s, y) \in R \}| \leq b(\min s)$ and for every~$(s, y) \in R$, $y \geq \min s$. Indeed, given a function $g \in \P_b$, one can define the relation $R_g \in \R_b$ defined by $R_g = \{ (s, y) : y \in g(s) \}$ and given a relation $R \in \R_b$, the function $g_R \in \P_b$ is defined by $g_R(s) = \{ y : (s, y) \in R \}$.

Note that the map $g \mapsto R_g$ is computable, but the map $R \mapsto g_R$ is not even continuous. Thankfully, there is a Turing functional which, given $R$ and a finite set~$\rho$, decides whether $\rho$ is $g_R$-free or not. Indeed, to decide whether $\rho$ is $g_R$-free, one does not need to know $g_R$ restricted to $[\rho]^{\leq h(\cdot)}$, only to know $R$ restricted to $[\rho]^{\leq h(\cdot)} \times \rho$.

We are now ready to define the forcing question:

\begin{definition}
Let $(f, \sigma, X)$ be a condition and $\varphi(G)$ be a $\Sigma^0_1$-formula.
Let $(f, \sigma, X) \qvdash \varphi(G)$ hold if and only if for every relation $R \in \R_{b_f^{+}}$, there is a finite $g_R$-free set $\rho \subseteq X$ such that $\varphi(\sigma \cup \rho)$ holds.
\end{definition}

The previous formulation of the forcing question is $\Pi^1_1(X)$ as it starts with a second-order universal quantification. However, thanks to the effective compactness of the space $\R_{b_f^{+}}$, it is equivalent to a $\Sigma^0_1(X)$-formula:

\begin{lemma}\label[lemma]{lem:fs-forcing-question-sigma01}
Let $(f, \sigma, X)$ be a condition and $\varphi(G)$ be a $\Sigma^0_1$-formula.
Then $(f, \sigma, X) \qvdash \varphi(G)$ if and only if there is some~$k \in \omega$ such that for every $b_f^{+}$-constrained progressive function $g : [0, k]^{\leq h(\cdot)} \to [0, k]^{<\omega}$, there is a finite $g$-free set $\rho \subseteq X \uh_k$ such that $\varphi(\sigma \cup \rho)$ holds.
\end{lemma}
\begin{proof}
Suppose first that there is some~$k \in \NN$ such that for every $b_f^{+}$-constrained progressive function $g : [0, k]^{\leq h(\cdot)} \to [0, k]^{<\omega}$, there is a finite $g$-free set $\rho \subseteq X \uh_k$ such that $\varphi(\sigma \cup \rho)$ holds. Let $R \in \R_{b_f^+}$ be a relation,
and let $g_R \in \P_{b_f^+}$ be the corresponding function. Then, letting $g : [0, k]^{\leq h(\cdot)} \to [0, k]^{<\omega}$ be defined by $g(s) = g_R(s) \cap [0, k]$, the function $g$ is $b_f^{+}$-constrained and progressive, so there is some finite $g$-free set $\rho \subseteq X \uh_k$ such that $\varphi(\sigma \cup \rho)$ holds. In particular, $\rho$ is $g_R$-free. Since there is such a $\rho$ for every $R \in \R_{b_f^+}$, then $(f, \sigma, X) \qvdash \varphi(G)$ holds.

Suppose now that for every~$k \in \NN$, there is a $b_f^{+}$-constrained progressive function $g_k : [0, k]^{\leq h(\cdot)} \to [0, k]^{<\omega}$ such that for every $g_k$-free set $\rho \subseteq X \uh_k$, $\varphi(\sigma \cup \rho)$ does not hold. Let $\T$ be the tree which, at level~$k$, contains all such functions $g_k$, and which is ordered by the function extension relation. The tree $\T$ is finitely branching, so by K\"onig's lemma, there is an infinite path $g \in \T$. This path is a function $g \in \P_{b_f^+}$ such that for every finite $g$-free set $\rho \subseteq X$, $\varphi(\sigma \cup \rho)$ does not hold. Then the relation $R_g$ witnesses that $(f, \sigma, X) \qvdash \varphi(G)$ does not hold.
\end{proof}

Any filter $\F$ for this notion of forcing induces a set 
$$
G_\F = \bigcup \{ \sigma : (g, \sigma, X) \in \F \}
$$
We say that a condition $p$ \emph{forces} a formula $\varphi(G)$ if $\varphi(G_\F)$ holds for every sufficiently generic filter~$\F$ containing~$p$.
The following lemma states that the forcing question meets its specification.

\begin{lemma}\label[lemma]{lem:fs-forcing-question-spec}
Let $p = (f, \sigma, X)$ be a condition and $\varphi(G)$ be a $\Sigma^0_1$-formula.
\begin{enumerate}
    \item[1.] If $p \qvdash \varphi(G)$, then there is an extension of~$p$ forcing $\varphi(G)$.
    \item[2.] If $p \nqvdash \varphi(G)$, then there is an extension of~$p$ forcing $\neg \varphi(G)$.
\end{enumerate}
\end{lemma}
\begin{proof}
Suppose first $p \qvdash \varphi(G)$ holds. By \Cref{lem:fs-forcing-question-sigma01}, there is some~$k \in \NN$ such that for every $b_f^{+}$-constrained progressive function $g : [0, k]^{\leq h(\cdot)} \to [0, k]^{<\omega}$, there is a finite $g$-free set $\rho \subseteq X$ such that $\varphi(\sigma \cup \rho)$ holds. By \Cref{lem:fs-stabilizing}, there is an infinite subset $Y_0 \subseteq X$ stabilizing $[0, k]$ and such that $D \not \leq_T Y_0$. Let $g_{k,Y_0} : [0, k]^{\leq h(\cdot)} \to [0, k]^{<\omega}$ be the limit function. Note that $g_{k,Y_0}$ is $b_f^{+}$-constrained and progressive, so there is a finite $g_{k,Y_0}$-free set $\rho \subseteq X$ such that $\varphi(\sigma \cup \rho)$ holds.
By \Cref{lem:fs-preserving-b}, $(f, \sigma \cup \rho, Y_0)$ satisfies (b).
By \Cref{lem:fs-preserving-a}, there is an infinite subset $Y \subseteq Y_0$ such that $(f, \sigma \cup \rho, Y)$ satisfies (a) and $D \not \leq_T Y$. Thus $(f, \sigma \cup \rho, Y)$ is a valid extension of~$p$. By choice of~$\rho$, it forces $\varphi(G)$.

Suppose $p \nqvdash \varphi(G)$ holds. Then there is a relation $R \in \R_{b_f^{+}}$ such that for every finite $g_R$-free set $\rho \subseteq X$, $\varphi(\sigma \cup \rho)$ does not hold. Let $g_R \in \P_{b_f^{+}}$ be the corresponding function. Let $\hat{f} : [\NN]^{\leq h(\cdot)} \to [\NN]^{<\omega}$ be defined by $\hat{f}(s) = f(s) \cup (g_R(s) \setminus \sigma)$. Note that $\hat{f} \leq f$, and every $\hat{f}$-free subset $\rho \subseteq X$ is $g_R$-free. Moreover, $\hat{f}$ is $(b_f + b_f^{+})$-constrained and progressive. Since $(f, \sigma, X)$ satisfies (b) and for every $s \in [\sigma \cup X]^{\leq h(\cdot)}$, $\hat{f}(s)\cap \sigma\subseteq f(s)\cap \sigma\subseteq s$, then $(\hat{f}, \sigma, X)$ satisfies (b). By \Cref{lem:fs-preserving-a}, there is an infinite subset $Y \subseteq X$ such that $D \not \leq_T Y$ and $(\hat{f}, \sigma, Y)$ satisfies (a). Thus $q = (\hat{f}, \sigma, Y)$ is a valid extension of~$p$. 

We claim that $q$ forces $\neg \varphi(G)$. Indeed, suppose there is an extension $(\tilde{f}, \tau, Z) \leq q$ such that $\varphi(\tau)$ holds. Then by property (b) of  $(\tilde{f}, \tau, Z)$, $\tau$ is $\tilde{f}$-free, and since $\tilde{f} \leq \hat{f}$, $\tau$ is $\hat{f}$-free. Let $\rho = \tau \setminus \sigma$. By definition of~$\hat{f}$, $\rho$ is a $g_R$-free subset of~$X$, contradicting our choice of~$g_R$.
\end{proof}

We can now prove our diagonalization lemma.

\begin{lemma}\label[lemma]{lem:fs-diagonalization}
Let $p = (f, \sigma, X)$ be a condition and $\Phi_e$ be a Turing functional.
There is an extension of~$p$ forcing $\Phi_e^G \neq D$.
\end{lemma}
\begin{proof}
Let $U = \{ (x, v) \in \NN \times 2 : p \qvdash \Phi^G_e(x) \downarrow = v \}$. We have three cases:
\begin{itemize}
    \item Case 1: $(x, 1-D(x)) \in U$ for some~$x \in \NN$. By \Cref{lem:fs-forcing-question-spec}, there is an extension of~$p$ forcing $\Phi^G_e(x)\downarrow = 1-D(x)$, hence forcing $\Phi_e^G \neq D$.
    \item Case 2: $(x, D(x)) \not \in U$ for some~$x \in \NN$. By \Cref{lem:fs-forcing-question-spec}, there is an extension of~$p$ forcing $\neg (\Phi^G_e(x)\downarrow = D(x))$, hence forcing $\Phi_e^G \neq D$.
    \item Case 3: $U$ is the graph of the characteristic function of~$D$. By \Cref{lem:fs-forcing-question-sigma01}, the set $U$ is $\Sigma^0_1(X)$, so $D \leq_T X$, contradiction.
\end{itemize}
\end{proof}

We are now ready to prove \Cref{thm:pfs-omega-strong-cone-avoidance}.
Let $f : [\NN]^{\leq h(\cdot)} \to [\NN]^{<\omega}$ be an $h$-constrained, progressive coloring, for a computable function $h : \NN \to \NN$.
Let $\F$ be a sufficiently generic filter containing $(f, \emptyset, \NN)$,
and let $G_\F = \bigcup \{ \sigma : (g, \sigma, X) \in \F \}$. By definition of a forcing condition, $G_\F$ is $f$-free. By \Cref{lem:fs-extendibility}, $G_\F$ is infinite, and by \Cref{lem:fs-diagonalization}, $D \not \leq_T G_\F$.
This completes the proof of \Cref{thm:pfs-omega-strong-cone-avoidance}.
\end{proof}

\bibliographystyle{plain}
\bibliography{biblio}
\end{document}
